\newpage
\phantomsection
\label{prf:fraction-order-comparison-general}

\begin{remark*}[Return]
\hyperref[thm:fraction-order-comparison-general]{Return to Theorem}
\end{remark*}

\begin{theorem*}[General Fraction Comparison in an Ordered Field]
For any $a,b,c,d\in F$ with $b\neq 0$ and $d\neq 0$,
\[
\frac{a}{b}<\frac{c}{d}
\quad\Longleftrightarrow\quad
(a\cdot d)\cdot(b\cdot d)<(b\cdot c)\cdot(b\cdot d).
\]
\end{theorem*}

\begin{proof}
\textbf{Professional Standard Proof.}~\\
Let $a,b,c,d\in F$ with $b\neq 0$ and $d\neq 0$.  Since $F$ is a field,
\[
b\cdot d\neq 0.
\]

We begin by rewriting the comparison of the two quotients as a positivity
statement about their difference:
\[
\frac{a}{b}<\frac{c}{d}
\quad\Longleftrightarrow\quad
0<\frac{c}{d}-\frac{a}{b}.
\]

Set
\[
x:=\frac{c}{d}-\frac{a}{b}.
\]
We claim that
\[
x=\frac{b\cdot c-a\cdot d}{b\cdot d}.
\]
To verify this, multiply $x$ by $b\cdot d$:
\[
(b\cdot d)x
=(b\cdot d)\left(\frac{c}{d}-\frac{a}{b}\right).
\]
Distributing and regrouping,
\[
(b\cdot d)x
=(b\cdot d)\frac{c}{d}-(b\cdot d)\frac{a}{b}
=b\left(d\frac{c}{d}\right)-d\left(b\frac{a}{b}\right).
\]
By the defining property of division,
\[
d\frac{c}{d}=c
\qquad\text{and}\qquad
b\frac{a}{b}=a.
\]
Hence
\[
(b\cdot d)x=b\cdot c-a\cdot d.
\]
By the definition of division, the unique element whose product with
$b\cdot d$ is $b\cdot c-a\cdot d$ is
\[
\frac{b\cdot c-a\cdot d}{b\cdot d}.
\]
Therefore
\[
\frac{c}{d}-\frac{a}{b}=\frac{b\cdot c-a\cdot d}{b\cdot d}.
\]

Substituting this into the inequality above gives
\[
\frac{a}{b}<\frac{c}{d}
\quad\Longleftrightarrow\quad
0<\frac{b\cdot c-a\cdot d}{b\cdot d}.
\]
Now apply
\hyperref[thm:positive-fraction-iff-positive-product]{Theorem~\ref*{thm:positive-fraction-iff-positive-product}}
with numerator $b\cdot c-a\cdot d$ and denominator $b\cdot d$:
\[
0<\frac{b\cdot c-a\cdot d}{b\cdot d}
\quad\Longleftrightarrow\quad
0<(b\cdot c-a\cdot d)\cdot(b\cdot d).
\]
Distributing the product on the right,
\[
0<(b\cdot c)\cdot(b\cdot d)-(a\cdot d)\cdot(b\cdot d).
\]
Adding $(a\cdot d)\cdot(b\cdot d)$ to both sides, we obtain
\[
(a\cdot d)\cdot(b\cdot d)<(b\cdot c)\cdot(b\cdot d).
\]

Combining the equivalences yields
\[
\frac{a}{b}<\frac{c}{d}
\quad\Longleftrightarrow\quad
(a\cdot d)\cdot(b\cdot d)<(b\cdot c)\cdot(b\cdot d),
\]
as required.
\end{proof}

\begin{remark*}[Proof structure]
The proof converts the comparison of two quotients into a positivity statement
for their difference.  That difference is then rewritten as a single quotient
with denominator $bd$ by checking directly that multiplying by $bd$ recovers
the numerator $bc-ad$.  Once the expression has been packed into one quotient,
the earlier theorem characterizing positivity of a quotient by the sign of its
numerator-times-denominator applies immediately.  A final additive
rearrangement turns the positivity condition into the comparison
$(ad)(bd)<(bc)(bd)$.
\end{remark*}

\begin{remark*}[Dependencies]~\\
\begin{itemize}
  \item \hyperref[def:ordered-field]{Definition~\ref*{def:ordered-field}}
  \item \hyperref[thm:positive-fraction-iff-positive-product]{Theorem~\ref*{thm:positive-fraction-iff-positive-product}}
\end{itemize}
\end{remark*}
